\section{Conclusiones y Líneas Futuras}
\label{sec:conclusiones}

El análisis realizado sobre la prima de riesgo griega revela la dualidad de los mercados financieros en periodos de crisis. Por un lado, el enfoque de la \textbf{Econofísica} ha permitido identificar una estructura de memoria a largo plazo ($\alpha \approx 0.94$), sugiriendo que la volatilidad no sigue un paseo aleatorio, sino que responde a dinámicas persistentes y comportamientos de rebaño.

Por otro lado, el modelo \textbf{GARCH Bayesiano}, a pesar de la complejidad computacional y las dificultades de convergencia observadas en las cadenas de Markov, ha demostrado ser una herramienta superior para la monitorización del riesgo a corto plazo, capturando los clústeres de volatilidad reaccionando a la información diaria.

\textbf{Trabajo Futuro:}
Para investigaciones posteriores, se propone extender la calibración Bayesiana utilizando servidores de alto rendimiento que permitan aumentar el número de muestras (draws) y cadenas, garantizando así un estadístico $\hat{R} \approx 1.0$. Asimismo, sería relevante aplicar este marco comparativo a otros países periféricos de la Eurozona para confirmar si este comportamiento super-difusivo es sistémico o idiosincrásico de Grecia.